\subsection{Księga pierwsza -- pojęcia podstawowe}
\setcounter{counter_euclid}{0}
% Pojęcie podstawowe 1
\begin{euclid_basic}
Wyrażenia, które są równe się temu samemu wyrażeniowi, są sobie równe.
\end{euclid_basic}

% Pojęcie podstawowe 2
\begin{euclid_basic}
Jeżeli równania dodawane są do równań, wtedy całości są sobie równe.
\end{euclid_basic}

% Pojęcie podstawowe 3
\begin{euclid_basic}
Jeżeli równania odejmowane są do równań, wtedy całości są sobie równe.
\end{euclid_basic}

% Pojęcie podstawowe 4
\begin{euclid_basic}
Wyrażenia, które się pokrywają, są sobie równe.
\end{euclid_basic}

% Pojęcie podstawowe 5
\begin{euclid_basic}
Całość jest większa od części. 
\end{euclid_basic}
